\documentclass[12pt,a4paper,reqno]{amsart}

% language
\usepackage[greek,english]{babel}
\usepackage[utf8]{inputenc}
\usepackage{alphabeta}

% change default names to greek
\addto\captionsenglish{
    \renewcommand{\contentsname}{Περιεχόμενα}
    \renewcommand{\refname}{Βιβλιογραφία}
    \renewcommand{\datename}{Ημερομηνία:}
    \renewcommand{\urladdrname}{Ιστοσελίδα}
}

% math 
\usepackage{amsmath,amsthm,amssymb,amscd,stmaryrd} 

% font
\usepackage[cal=euler]{mathalfa}
\usepackage{libertinus-type1}
% \usepackage{txfonts} % for upright greek letters
\usepackage{bm} % for bold symbols
\usepackage{bbm} % for the simply-looking bb symbols

% miscellaneous 
\usepackage{changepage} %for indenting environments
\usepackage{csquotes} % example: \textcquote{}
\usepackage{blkarray}
\setcounter{MaxMatrixCols}{20} % default for pmatrix is 10!!
\usepackage{ytableau}
\usepackage{array} %needed to increase the vertical length in a tabular

% drawing
\usepackage{tikz,tikz-cd}
\usetikzlibrary{shapes.misc, patterns, matrix, calc, intersections,positioning,arrows.meta}
\usepackage{graphics,graphicx}
\usepackage{float} % provides enhanced control and customization options for floating objects such as figures and tables

% colors
\usepackage{xcolor}
\definecolor{darkcandyapplered}{rgb}{0.64, 0.0, 0.0}
\definecolor{midnightblue}{rgb}{0.1, 0.1, 0.44}
\definecolor{mylightblue}{HTML}{336699}
\definecolor{burntorange}{rgb}{0.8, 0.33, 0.0}
\definecolor{iceberg}{rgb}{0.44, 0.65, 0.82}
\definecolor{applegreen}{rgb}{0.55, 0.71, 0.0}
\definecolor{canaryyellow}{rgb}{1.0, 0.94, 0.0}
\definecolor{lightgray}{rgb}{0.83, 0.83, 0.83}

% hrefs
\usepackage{hyperref}
\usepackage[noabbrev,capitalize]{cleveref}
\hypersetup{
    pdftoolbar=true,        
    pdfmenubar=true,        
    pdffitwindow=false,     
    pdfstartview={FitH},  % fits the width of the page to the window
    pdftitle={},
    pdfauthor={},
    pdfsubject={},
    pdfkeywords={},
    pdfnewwindow=true,  % links in new window
    colorlinks=true,  % false: boxed links; true: colored links
    linkcolor=darkcandyapplered,   % color of internal links
    citecolor=midnightblue,  % color of links to bibliography
    urlcolor=cyan,  % color of external links
    linktocpage=true  % changes the links from the section body to the page number
    }

% geometry
\textwidth=16cm 
\textheight=21cm 
\hoffset=-55pt 
\footskip=25pt

% thm envs (you might need to change the path)
\theoremstyle{definition}
\newtheorem{theorem}{Θεώρημα}
\newtheorem{proposition}[theorem]{Πρόταση}
\newtheorem{lemma}[theorem]{Λήμμα}
\newtheorem{corollary}[theorem]{Πόρισμα}
\newtheorem{conjecture}[theorem]{Εικασία}
\newtheorem{problem}[theorem]{Πρόβλημα}
\newtheorem*{claim}{Ισχυρισμός}
\newtheorem{observation}[theorem]{Παρατήρηση}
\newtheorem{definition}[theorem]{Ορισμός}
\newtheorem{question}[theorem]{Ερώτημα}
\newtheorem*{questions}{Ερωτήματα}
\newtheorem*{que}{Ερώτημα}
\newtheorem{example}[theorem]{Παράδειγμα}
\newtheorem{exercise}{Άσκηση}
\newtheorem*{zorn}{Λήμμα του Zorn}
\newtheorem*{example_6.6}{Παράδειγμα~6.6}

\theoremstyle{remark}
\newtheorem*{remark}{Παρατήρηση}

% fixes the correct numbering of environments
\numberwithin{theorem}{section}
\numberwithin{exercise}{section}
\numberwithin{equation}{section}

% math ops 
% fields
\renewcommand{\AA}{\mathbbmss{A}} 
\newcommand{\NN}{\mathbbmss{N}} 
\newcommand{\ZZ}{\mathbbmss{Z}} 
\newcommand{\QQ}{\mathbbmss{Q}} 
\newcommand{\RR}{\mathbbmss{R}} 
\newcommand{\CC}{\mathbbmss{C}} 
\newcommand{\KK}{\mathbbmss{K}} 
\newcommand{\FF}{\mathbbmss{F}} 
\newcommand{\HH}{\mathbbmss{H}} 

% frak
\newcommand{\fS}{\mathfrak{S}}  
\newcommand{\fm}{\mathfrak{m}}  
\newcommand{\fp}{\mathfrak{p}}  

% calligraphic 
\newcommand{\aA}{\mathcal{A}} 
\newcommand{\bB}{\mathcal{B}}
\newcommand{\cC}{\mathcal{C}}
\newcommand{\dD}{\mathcal{D}}
\newcommand{\eE}{\mathcal{E}}
\newcommand{\fF}{\mathcal{F}}
\newcommand{\hH}{\mathcal{H}}
\newcommand{\iI}{\mathcal{I}}
\newcommand{\lL}{\mathcal{L}}
\newcommand{\oO}{\mathcal{O}}
\newcommand{\pP}{\mathcal{P}}
\newcommand{\qQ}{\mathcal{Q}}
\newcommand{\sS}{\mathcal{S}}
\newcommand{\mM}{\mathcal{M}}
\newcommand{\uU}{\mathcal{U}}
\newcommand{\vV}{\mathcal{V}}
\newcommand{\zZ}{\mathcal{Z}}

% bold
\newcommand{\bfa}{\pmb{a}}
\newcommand{\bfe}{\mathbf{e}}
\newcommand{\bfF}{\pmb{F}}
\newcommand{\bfR}{\pmb{R}}
\newcommand{\bfv}{\mathbf{v}}
%\newcommand{\bfx}{\bm{x}}
%\newcommand{\bfx}{\mathbf{x}} 
\newcommand{\bfx}{\pmb{x}}
\newcommand{\bfX}{\pmb{X}}
\newcommand{\bfy}{\pmb{y}}
\newcommand{\bfz}{\pmb{z}}

% roman
\newcommand{\rmA}{\mathrm{A}}
\newcommand{\rmB}{\mathrm{B}}
\newcommand{\rmC}{\mathrm{C}}
\newcommand{\rmD}{\mathrm{D}} 
\newcommand{\rmF}{\mathrm{F}} 
\newcommand{\rmH}{\mathrm{H}} 
\newcommand{\rmh}{\mathrm{h}} 
\newcommand{\rmI}{\mathrm{I}} 
\newcommand{\rmK}{\mathrm{K}}
\newcommand{\rmM}{\mathrm{M}}
\newcommand{\rmP}{\mathrm{P}}  
\newcommand{\rmp}{\mathrm{p}}  
\newcommand{\rmQ}{\mathrm{Q}}  
\newcommand{\rmR}{\mathrm{R}}
\newcommand{\rmS}{\mathrm{S}}
\newcommand{\rmT}{\mathrm{T}}
\newcommand{\rmU}{\mathrm{U}}
\newcommand{\rmV}{\mathrm{V}}
\newcommand{\rmY}{\mathrm{Y}}
\newcommand{\rmZ}{\mathrm{Z}}
\newcommand{\rmz}{\mathrm{z}}

% arrows and symbols 
\renewcommand{\to}{\rightarrow}
\newcommand{\toto}{\longrightarrow}
\newcommand{\mapstoto}{\longmapsto}
\newcommand{\then}{\Rightarrow}
\newcommand{\IFF}{\Leftrightarrow}
\newcommand{\tl}{\tilde}
\newcommand{\wtl}{\widetilde}
\newcommand{\ol}{\overline}
\newcommand{\ul}{\underline}
\newcommand{\oldemptyset}{\emptyset}
\renewcommand{\emptyset}{\varnothing}
\DeclareMathSymbol{\Arg}{\mathbin}{AMSa}{"39} % for arguments 
\newcommand{\onto}{\ensuremath{\twoheadrightarrow}}
\newcommand{\tle}{\trianglelefteq}
\newcommand{\tge}{\trianglerighteq}

% custom ops
\newcommand{\rmKer}{\mathrm{Ker}}
\newcommand{\rmIm}{\mathrm{Im}}
\newcommand{\lcm}{\mathrm{lcm}}
\newcommand{\GL}{\mathrm{GL}}
\newcommand{\ev}{\mathrm{ev}}
\newcommand{\End}{\mathrm{End}}
\newcommand{\rmid}{\mathrm{id}}
\newcommand{\rmspan}{\mathrm{span}}
\newcommand{\Hom}{\mathrm{Hom}}
\newcommand{\Ann}{\mathrm{Ann}}
\newcommand{\Set}{\pmb{\mathrm{Set}}}
\newcommand{\Rmod}{\pmb{\mathrm{mod}}}
\newcommand{\rmrank}{\mathrm{rank}}
\newcommand{\mSpec}{\mathrm{mSpec}}
\newcommand{\Spec}{\mathrm{Spec}}
\newcommand{\tor}{\mathrm{tor}}
\newcommand{\rad}{\mathrm{rad}}
\newcommand{\Pol}[1]{F[x_1, x_2, \dots, x_{#1}]}

% absolute value symbol
\usepackage{mathtools} 
\DeclarePairedDelimiter\abs{\lvert}{\rvert}%
\DeclarePairedDelimiter\norm{\lVert}{\rVert}%
\makeatletter
\let\oldabs\abs
\def\abs{\@ifstar{\oldabs}{\oldabs*}}

% tensor symbol
\newcommand{\tensor}[1]{%
  \mathbin{\mathop{\otimes}\limits_{#1}}%
}

% permutation cycle notation
\ExplSyntaxOn
\NewDocumentCommand{\cycle}{ O{\;} m }
 {
  (
  \alec_cycle:nn { #1 } { #2 }
  )
 }

\seq_new:N \l_alec_cycle_seq
\cs_new_protected:Npn \alec_cycle:nn #1 #2
 {
  \seq_set_split:Nnn \l_alec_cycle_seq { , } { #2 }
  \seq_use:Nn \l_alec_cycle_seq { #1 }
 }
\ExplSyntaxOff

% setminus symbol
\newcommand{\mysetminusD}{\hbox{\tikz{\draw[line width=0.6pt,line cap=round] (3pt,0) -- (0,6pt);}}}
\newcommand{\mysetminusT}{\mysetminusD}
\newcommand{\mysetminusS}{\hbox{\tikz{\draw[line width=0.45pt,line cap=round] (2pt,0) -- (0,4pt);}}}
\newcommand{\mysetminusSS}{\hbox{\tikz{\draw[line width=0.4pt,line cap=round] (1.5pt,0) -- (0,3pt);}}}
\newcommand{\sm}{\mathbin{\mathchoice{\mysetminusD}{\mysetminusT}{\mysetminusS}{\mysetminusSS}}}

% miscellaneous commands
\newcommand{\defn}[1]{{\color{mylightblue}{#1}}}
\newcommand{\toDo}{{\bf\color{red} TODO}}
\newcommand{\toCite}{{\bf\color{green} CITE}}
\newcommand*{\vertbar}{\rule[-1ex]{0.5pt}{2.5ex}} % for matrices with column vectors
\newcommand*{\horzbar}{\rule[.5ex]{2.5ex}{0.5pt}} % for matrices with row vectors
\newcommand{\myblue}[1]{{\color{iceberg}{#1}}}
\newcommand{\myorange}[1]{{\color{burntorange}{#1}}}
\newcommand{\mygreen}[1]{{\color{applegreen}{#1}}}
\newcommand{\myred}[1]{{\color{darkcandyapplered}{#1}}}

% ferrer's diagram
\newcommand{\fdiagram}[1]{
    \begin{tikzpicture}[scale=.7]
        \fill foreach \Z [count=\Y] in {#1}
        {foreach \X in {1,...,\Z} 
        {(\X,-\Y) circle[radius=3pt]}};
    \end{tikzpicture}
}

%
\usepackage{pgfplots}
\usepgfplotslibrary{groupplots}
\pgfplotsset{compat=1.18}

%
\makeatletter
\def\mathcolor#1#{\@mathcolor{#1}}
\def\@mathcolor#1#2#3{%
  \protect\leavevmode
  \begingroup
    \color#1{#2}#3%
  \endgroup
}
\makeatother

% restriction
\newcommand\restr[2]{{% we make the whole thing an ordinary symbol
  \left.\kern-\nulldelimiterspace % automatically resize the bar with \right
  #1 % the function
  \littletaller % pretend it's a little taller at normal size
  \right|_{#2} % this is the delimiter
  }}
\newcommand{\littletaller}{\mathchoice{\vphantom{\big|}}{}{}{}}

% 
\newenvironment{nouppercase}{%
  \let\uppercase\relax%
  \renewcommand{\uppercasenonmath}[1]{}}{}

% titlepage
\title{226: Θεωρία Δακτυλίων και Modules}
\author[Β.~Δ. Μουστακας]{Βασίλης Διονύσης Μουστάκας \\ Πανεπιστήμιο Κρήτης}
\date{14 Μαΐου 2026}
% \urladdr{\href{https://sites.google.com/view/vasmous}{https://sites.google.com/view/vasmous}}

\begin{document}

\begingroup
\def\uppercasenonmath#1{} % this disables uppercase title
\let\MakeUppercase\relax % this disables uppercase authors
\maketitle
\endgroup

\setcounter{section}{11}
\setcounter{theorem}{0}
\setcounter{equation}{0}
\begin{center}
    \textbf{11. Το Θεώρημα μηδενικών του Hilbert
} 
\end{center}

Στην Ενότητα 9, κατασκευάσαμε απεικονίσεις  
\begin{align*}
\left\{
\begin{aligned}
&\qquad \quad \text{ιδεώδη} \\
& I \subseteq \Pol{n}
\end{aligned}
\right\}
&\to
\left\{
\begin{aligned}
&\text{πολλαπλότητες} \\
&\quad \ \ V \subseteq \AA_F^n
\end{aligned}
\right\} \\
I &\mapsto \vV(I) \\
\iI(V) &\mapsfrom V,
\end{align*}
όπου $\vV(I)$ είναι η \emph{πολλαπλότητα} που αντιστοιχεί στο $I$ και $\iI(A)$ είναι το \emph{ιδεώδες μηδενισμού} που αντιστοιχεί στο $A$. Αυτές ικανοποιούν τις σχέσεις 
\begin{align*}
    V &= \vV\left(\iI(V)\right) \\
    I &\subseteq \iI\left(\vV(I)\right)
\end{align*}
(βλ. Πρόταση 9.4). 

Από την μία μεριά, στο Παράδειγμα 9.2 (ζ) είδαμε ότι η απεικόνιση $\vV(\cdot)$ δεν είναι ένα προς ένα. Από την άλλη μεριά, η Πρόταση 9.10 μας πληροφορεί ότι υπάρχουν ιδεώδη τα οποία δεν είναι ιδεώδη μηδενισμού κάποιας πολλαπλότητας\footnote{Αρκεί να διαλέξουμε οποιοδήποτε μη-ριζικό ιδεώδες, όπως το $(x^2)$ στο $F[x]$.}. Με άλλα λόγια, απεικόνιση $\iI(\cdot)$ δεν είναι επί. Πότε είναι η μία αντίστροφη της άλλης; Το επόμενο αποτέλεσμα δίνει μια απάντηση στο ερώτημα αυτό.

\begin{theorem}
  \label{thm:mSpec}
  Αν το $F$ είναι αλγεβρικά κλειστό, τότε κάθε μεγιστικό ιδεώδες του $F[x_1, x_2,$ $\dots, x_n]$ είναι της μορφής $\fm_{\bfa}$, για κάποιο $\bfa \in \AA^n$. Ειδικότερα, υπάρχει μια αμφιμονοσήμαντη αντιστοιχία 
  \[
  \mSpec\left(
    \Pol{n}
  \right) \to \AA^n.
  \]
\end{theorem}

\begin{remark}
  Το Θεώρημα \ref{thm:mSpec} παύει να ισχύει όταν το $F$ δεν είναι αλγεβρικά κλειστό. Για παράδειγμα, για $F = \RR$, το ιδεώδες που παράγεται από το πολυώνυμο $1 + x^2$ είναι μεγιστικό, διότι όπως είδαμε στην Ενότητα 3, $\RR[x]/(1+x^2) \cong \CC$, το οποίο είναι σώμα, αλλά δεν αντιστοιχεί σε κάποιο σημείο του $\AA_\RR^1$ όπως είδαμε στο Παράδειγμα 9.2 (η).
\end{remark}

Ας συζητήσουμε μερικές σημαντικές συνέπειες του Θεωρήματος \ref{thm:mSpec}, την απόδειξη του οποίου αναβάλλουμε για την επόμενη ενότητα. Στο παράδειγμα 9.2 (η) είδαμε ότι 
\[
\vV(1 + x^2) = \vV(1 + x^2 + x^4) = \emptyset \ \subset \ \AA_\RR^1.
\]
Με άλλα λόγια, το σύστημα πολυωνυμικών εξισώσεων 
\begin{align*}
  1 + x^2 &= 0 \\
  1 + x^2 + x^4&= 0 
\end{align*}
δεν έχει λύση στο $\AA_\RR^1$. Πότε έχει ένα σύστημα πολυωνυμικών εξισώσεων κοινή λύση;

Επειδή το $F[x]$ είναι περιοχή κύριων ιδεωδών (βλ. συζήτηση στην αρχή της Ενότητας 3), κάθε ιδεώδες είναι της μορφής $(f)$ για κάποιο $f \in F[x]$. Στο Παράδειγμα 9.2 (γ) είδαμε ότι η αντίστοιχη πολλαπλότητα $\vV(f)$ αποτελείται από όλες τις ρίζες του $f$. Αν το $F$ είναι αλγεβρικά κλειστό, τότε κάθε μη σταθερό πολυώνυμο έχει ρίζα. Κατά συνέπεια, ο μόνος τρόπος να έχουμε $\vV(f) = \emptyset$ είναι όταν το $f \in F\sm\{0\}$. Με άλλα λόγια, αν το $F$ είναι αλγεβρικά κλειστό το μοναδικό ιδεώδες που αντιστοιχεί στην \textquote{κενή} πολλαπλότητα είναι ολόκληρος ο πολυωνυμικός δακτύλιος.

Το επόμενο θεώρημα, γνωστό ως το ασθενές θεώρημα μηδενικών (nullstellensatz\footnote{Γερμανική λέξη που είναι συνκαιρασμός των λέξεων null (μηδέν), stellen (θέσεις) και satz (θεώρημα).}) του Hilbert, μας πληροφορεί ότι (παραδόξως) το ίδιο ισχύει και στον πολυωνυμικό δακτύλιο πολλών μεταβλητών. 
\begin{theorem}{\rm(Hilbert's weak nullstellensatz)}
  \label{thm:weak_nullstellensatz}
  Αν το $F$ είναι αλγεβρικά κλειστό σώμα και $I$ είναι ένα γνήσιο ιδεώδες του $F[x_1, x_2, \dots, x_n]$, τότε $\vV(I) \neq \emptyset$.
\end{theorem}

\begin{proof}[Απόδειξη]
  Επειδή το $I$ είναι γνήσιο ιδεώδες του $\Pol{n}$ έπεται ότι υπάρχει κάποιο μεγιστικό ιδεώδες $\fm$ τ' οποίο το περιέχει (γιατί;). Από το Θεώρημα \ref{thm:mSpec} $\fm = \fm_{\bfa}$, για κάποιο $\bfa \in \AA_F^n$. Κατά συνέπεια, από την παρατήρηση μετά τον Ορισμό 9.1 έπεται ότι 
  \[
  \vV(\fm_{\bfa}) = \{\bfa\} \subseteq \vV(I)
  \]
  όπου η πρώτη ισότητα έπεται από το Παράδειγμα 9.2 (δ) και γι αυτό $\vV(I) \neq \emptyset$.
\end{proof}

\begin{remark}
  Όπως έχουμε ήδη δει, ένα σύστημα πολυωνυμικών εξισώσεων 
  \begin{align*}
    f_1(x_1, x_2, \dots, x_n) &= 0 \\
    f_2(x_1, x_2, \dots, x_n) &= 0 \\
    &\ \ \vdots \\
    f_m(x_1, x_2, \dots, x_n) &= 0  
  \end{align*}
  δεν έχει κοινή λύση στο $\AA_F^n$ αν και μόνο αν 
  \[
  \vV(\{f_1, f_2, \dots, f_m\}) = \emptyset.
  \]
  Αν το $F$ είναι αλγεβρικά κλειστό σώμα, τότε το Θεώρημα \ref{thm:weak_nullstellensatz} μας πληροφορεί ότι αυτό συμβαίνει αν και μόνο αν 
  \[
  1 \in \left(
    f_1, f_2, \dots, f_m
  \right),
  \]
  μετατρέποντας έτσι το αρχικό (γεωμετρικό) πρόβλημα στον καθορισμό του αν το $1$ ανήκει ή όχι σε ένα ιδεώδες.
\end{remark}

Στο τέλος της Ενότητας 9, μαντέψαμε ότι αυτό που εμποδίζει την απεικόνιση $I \mapsto \vV(I)$ από το να είναι ένα-προς-ένα είναι το κατά πόσο το $I$ είναι ριζικό ιδεώδες. Το Θεώρημα \ref{thm:weak_nullstellensatz} μας επιτρέπει το δείξουμε με μόνη προϋπόθεση ότι δουλεύουμε πάνω από αλγεβρικά κλειστό σώμα. Το αποτέλεσμα αυτό είναι γνωστό ως το \emph{ισχυρό θεώρημα μηδενικών του Hilbert}.
\begin{theorem}{\rm(Hilbert's strong nullstellensatz)}
  \label{thm:strong_nullstellensatz}
  Υποθέτουμε ότι το $F$ είναι αλγεβρικά κλει\-στό σώμα. Αν $I$ είναι ιδεώδες του $F[x_1, x_2, \dots, x_n]$, τότε 
  \[
  \iI\left(
    \vV\left(
      I
    \right)
  \right) = \sqrt{I}.
  \]
   Ειδικότερα, υπάρχει μια αμφιμονοσήμαντη αντιστοιχία 
  \begin{align*}
  \left\{
  \begin{aligned}
  &\quad \ \text{ριζικά ιδεώδη} \\
  & I \subseteq \Pol{n}
  \end{aligned}
  \right\}
  &\to
  \left\{
  \begin{aligned}
  &\text{πολλαπλότητες} \\
  &\quad \ \ V \subseteq \AA_F^n
  \end{aligned}
  \right\}.
  \end{align*}
\end{theorem}

\begin{proof}[Απόδειξη]
  Για λόγους διευκόλυνσης στον συμβολισμό, σε ότι ακολουθεί γράφουμε 
  \[
  \bfx = (x_1, x_2, \dots, x_n).
  \]
  Ο δεύτερος ισχυρισμός έπεται άμεσα από τον πρώτο, διότι 
  \[
  \iI\left(
    \vV\left(
      I
    \right)
  \right) = \sqrt{I} = I
  \]
  για κάθε ριζικό ιδεώδες $I$. Οι λεπτομέρειες αφήνονται ως άσκηση. Θα αποδείξουμε τον πρώτο ισχυρισμό.

  Η κατεύθενση \textquote{$\supseteq$} έπεται ως εξής: Αν $f \in \sqrt{I}$, τότε $f^k \in I$, για κάποιο $k \ge 1$. Συνεπώς, $f^k(\bfa) = 0$, για κάθε $\bfa \in \vV(I)$. Αυτό σημαίνει ότι $f(\bfa) = 0$, για κάθε $\bfa \in \vV(I)$ και γι αυτό $f \in \iI\left(
    \vV\left(
      I
    \right)
  \right)$.

  Για τον αντίστροφο εγκλεισμό\footnote{Στη βιβλιογραφία, αυτός ο τρόπος να αποδείξει κανείς το ισχυρό θεώρημα μηδενικών από το ασθενές ονομάζεται \textquote{τέχνασμα του Rabinowitsch}.}, έστω $f \in \iI\left(
    \vV\left(
      I
    \right)
  \right)$. Αν το $f$ είναι το μηδενικό πολυώνυμο, τότε $f \in I$, διότι το $I$ είναι ιδεώδες και το ζητούμενο έπεται. Υποθέτουμε λοιπόν, ότι $f \neq 0$. Από το Θεώρημα 10.1, έπεται ότι 
  \[
  I = \left(
    f_1, f_2, \dots, f_m
  \right),
  \]
  για κάποια $f_1, f_2, \dots, f_m \in F[\bfx]$. Θεωρούμε το ιδεώδες 
  \[
  \hat{I} = \left(
    f_1, f_2, \dots, f_m , 1 -yf
  \right) \ \subseteq \ F[\bfx, y],
  \]
  όπου θεωρούμε κάθε $f_i$ ως στοιχείο του $F[\bfx,y]$ με τον προφανή τρόπο.

  Ισχυριζόμαστε ότι $\vV(\hat{I}) = \emptyset$. Πράγματι, έστω $\hat{\bfa} = \left(\bfa, a_{n+1}\right) \in \AA^{n+1}$. Διακρίνουμε περιπτώσεις.
  \begin{itemize}
    \item Αν $\bfa \in \vV(I)$, τότε $f(\bfa) = 0$, διότι $f \in \iI\left(
    \vV\left(
      I
    \right)
  \right)$. Συνεπώς, 
  \[
  \ev_{\hat{\bfa}}(1-yf) = 1 - a_{n+1}f(\bfa) = 1 \neq 0
  \]
  και γι αυτό $\hat{\bfa} \notin \vV(\hat{I})$.
    \item Αν $\bfa \notin \vV(I)$, τότε
    \[
    0 \neq f_i(\bfa) = \ev_{\hat{\bfa}}(f_i)
    \]
    για κάποιο $1 \le i \le m$. Οπότε, $\hat{\bfa} \notin \vV(\hat{I})$.
  \end{itemize}

  Από το Θεώρημα \ref{thm:weak_nullstellensatz} έπεται ότι $1 \in \hat{I}$ και γι αυτό 
  \begin{equation}
    \label{eq:null_help}
    1 = \sum_{i=1}^m g_i(\bfx,y)f_i(\bfx) + g(\bfx,y)(1-yf(\bfx)),
  \end{equation}
  για κάποια $g_i, g \in F[x_1, \dots, x_n, y]$. Θεωρούμε\footnote{Έχει νόημα, διότι $f \neq 0$.} τον ομομορφισμό δακτυλίων 
  \begin{align*}
    \phi : F[\bfx,y] &\to F(\bfx) \\
    x_i &\mapsto x_i \\
    y &\mapstoto \frac{1}{f},
  \end{align*}
  όπου\footnote{Το $F(\bfx)$ είναι το σώμα πηλίκων του $F[\bfx]$ (βλ. σημείωση στην αρχή της απόδειξης του Θεωρήματος 3.6).} $F(\bfx)$ είναι το σώμα των ρητών συναρτήσεων σε $n$ μεταβλητές, πάνω από το $F$, δηλαδή στοιχείων της μορφής $p/q$, για κάποια $p, q \in F[\bfx]$ με $q \neq 0$.

  Εφαρμόζοντας τον $\phi$ στην Ταυτότητα \eqref{eq:null_help} παίρνουμε 
  \[
  1 = \sum_{i=1}^m g_i\left(\bfx, \frac{1}{f(\bfx)}\right)f_i(\bfx) + g\left(\bfx,\frac{1}{f(\bfx)}\right)\left(
    1 - \frac{1}{f(\bfx)}f(\bfx)
  \right)
  \]
  ή ισοδύναμα 
  \begin{equation}
    \label{eq:null_helpp}
    1 = \sum_{i=1}^m g_i\left(\bfx, \frac{1}{f(\bfx)}\right)f_i(\bfx).
  \end{equation}
  Το δεξί μέλος της Ταυτότητας \eqref{eq:null_helpp} είναι άθροισμα ρητών συναρτήσεων με παρονομαστές δυνάμεις του $f$. Διαλέγοντας την μεγαλύτερη δύναμη που εμφανίζεται, έστω $k \ge 1$, και πολλαπλασιάζοντας και τα δυο μέλη με $f^k$ έπεται ότι 
  \[
  f^k = \sum_{i=1}^m h_i f_i \ \in \ I,
  \]
  για κάποια $h_1, h_2, \dots, h_m \in F[\bfx]$. Άρα, $f \in \sqrt{I}$ και το ζητούμενο έπεται.
\end{proof}

Το Θεώρημα \ref{thm:strong_nullstellensatz} μας πληροφορεί ότι αν δύο ιδεώδη ορίζουν την ίδια πολλαπλότητα, τότε αυτά διαφέρουν κατά μία δύναμη. Πιο συγκεκριμένα, έχουμε το εξής.
\begin{corollary}
  \label{cor:strong_nullstellensatz}
  Υποθέτουμε ότι το $F$ είναι αλγεβρικά κλειστό σώμα. Αν $I, J$ είναι ιδεώδη του $\Pol{n}$ τέτοια ώστε $\vV(I) = \vV(J)$, τότε $f^k \in J$, για κάποιο $k \ge 1$, για κάθε $f \in I$.
\end{corollary}

% \begin{thebibliography}{99}
%     \bibitem{Alu09}
%     P.~Aluffi,
%     \textit{Algebra, Chapter 0},
%     Grad. Stud. in Math. {\bf104},
%     Amer. Math. Soc., Providence, RI, 2009.
%     \bibitem{DF04}
%     D. S. Dummit και R. M. Foote, 
%     \textit{Abstract algebra},
%     third edition, Wiley, 2004. 
% \end{thebibliography}
\end{document}